\documentclass[11pt]{article}
\usepackage[T1]{fontenc}
\usepackage[utf8]{inputenc}
\usepackage[margin=1in]{geometry}
\usepackage{amsmath,amssymb,amsthm,mathtools,mathrsfs}
\usepackage{array,booktabs,tabularx}
\usepackage{microtype}
\usepackage[colorlinks=true,linkcolor=blue,citecolor=blue,urlcolor=blue]{hyperref}
\hypersetup{pdftitle={An arithmetic source audit of the October 5 quasi-Riemann hypothesis argument},pdfauthor={Drafted for Edward Baker}}
\setlength{\emergencystretch}{3em}
\Urlmuskip=0mu plus 1mu
\numberwithin{equation}{section}
\newtheorem{theorem}{Theorem}[section]
\newtheorem{proposition}[theorem]{Proposition}
\newtheorem{lemma}[theorem]{Lemma}
\newtheorem{corollary}[theorem]{Corollary}
\theoremstyle{definition}
\newtheorem{assumption}[theorem]{Source interface}
\newtheorem{definition}[theorem]{Definition}
\theoremstyle{remark}
\newtheorem{remark}[theorem]{Remark}
\newcommand{\OO}{\mathcal O}
\newcommand{\E}{\mathcal E}
\newcommand{\N}{\mathrm N}
\newcommand{\ind}{\mathbf 1}
\newcommand{\norm}[1]{\left\|#1\right\|}
\newcommand{\sfstar}{\sideset{}{^*}\sum}
\newcommand{\ray}{\mathscr X}
\newcommand{\dd}{\,d}
\newcommand{\eps}{\epsilon}

\title{An arithmetic source audit of the October 5\\
quasi-Riemann hypothesis argument}
\author{Drafted for Edward Baker}
\date{12 October 2026}

\begin{document}
\maketitle

\begin{abstract}
We consolidate an internal source audit of the October 5, 2026
quasi-Riemann hypothesis manuscript of OpenAI. The audit follows its
special cubic Gauss-sum family through the initial Poisson reduction,
theta reflection, quadratic large sieve, cube inversion, and full
two-Poisson return. The central bookkeeping facts are a local
exceptional-prime reindexing, a signed preimage count
\(\tau(r)\ind_{f'\mid k'}\), an admissible divisor-shifted child, and a
strict decrease by the squared positive-gap ratio. We give the
resulting finite induction and the final deductions from the original
mean square to the fixed-character exponent \(11/12+\delta\), Mellin
nonvanishing, and the Dirichlet norm lift, including principal and
imprimitive characters. These deductions are conditional on the
explicit analytic interfaces stated here and the identified classical
imports. Finite computations diagnose algebra and normalization; they
do not certify the analytic theorems. No additional unresolved defect
was identified in this internal audit, but independent mathematical
validation of the source's claimed zero-free theorem remains outstanding.
\end{abstract}

\tableofcontents

\section{Scope and the assembled statement}\label{sec:scope}

The source of this investigation is \emph{The Quasi-Riemann Hypothesis},
dated October 5, 2026, by OpenAI \cite{Source}. Its main theorem asserts
nonvanishing of every finite-order Hecke \(L\)-function over
\(K=\mathbb Q(\sqrt{-3})\) in \(\Re s>11/12\), and consequently
the same statement for every Dirichlet \(L\)-function. The purpose
here is to make the audited dependencies and deductions reviewable.
The manuscript records an audit of that argument; it does not attribute
an independently verified new proof of quasi-RH to the author for whom
it was drafted.

Three levels of evidence must be kept distinct. The arithmetic
identities and parameter calculations below have explicit derivations.
The source's analytic transformations are stated at the precise
interfaces inspected in the accompanying notes. Cubic theta automorphy
and cusp coefficients, global reciprocity, the quadratic large sieve,
prime-ideal counting, and Hecke continuation remain classical imports.
The detailed original proofs of every such theorem have not been
independently re-proved in this investigation.

\begin{theorem}[Conditional assembly]\label{thm:assembly}
Fix a finite-order Hecke character \(\nu\) of \(K\) and a finite set
\(S\) containing the primes above \(2,3\) and the conductor primes of
\(\nu\). Assume the arithmetic, theta-reflection, smooth-weight,
Poisson-return interfaces stated in this
manuscript, with their uniformities and their indicated classical
inputs. Then for every \(0<\vartheta\le1/10\) and \(\epsilon>0\)
there is a finite integer \(k=k(\vartheta,\epsilon)\), independent
of the compact support interval \(I\), such that
\begin{equation}\label{eq:main-moment}
 \sum_{0<\N(u)\le D^{1+\vartheta}}|A_u(D)|^2
 \ll_{\nu,S,I,\vartheta,\epsilon}
 \norm{W}_{C^k(I)}^2D^{2+\vartheta+\epsilon},
 \qquad D\ge2.
\end{equation}
For every fixed smooth \(W\) and every \(\delta>0\),
\begin{equation}\label{eq:fixed-saving}
 A_1(D)\ll_{\nu,S,W,\delta}D^{11/12+\delta}.
\end{equation}
With the classical prime-ideal theorem and Hecke continuation, these
bounds imply that \(L_K(s,\nu)\) has no zeros in \(\Re s>11/12\).
Applying this conclusion to Dirichlet norm lifts gives the same
qualitative half-plane for every Dirichlet character.
\end{theorem}

The source interfaces in this statement are hypotheses of the
assembled deductions. The audit inspected their proofs and found no
additional unresolved defect under the stated imports, but that review
status is not a substitute for independent mathematical validation.
In particular, finite check counts are not hypotheses or proofs of
Theorem~\ref{thm:assembly}.

\subsection{Exact source version}

All source references use the physical lines of the October 5
LaTeX file at openai/math commit
\begin{center}
\texttt{fd4aeeb2ee4fc729c18d98444fed42fd0529eeeb}.
\end{center}
Its SHA-256, grouped solely for legibility, is
\begin{center}\small\ttfamily
d9a8f15aa770cf88 3d0eabd2b775fad6\\
94ce20b44cba7928 f5c0c9a6d8750d4d.
\end{center}
Concatenating the four groups gives the recorded hash. The commit
date is October 8; the manuscript date is October 5. They are
different pieces of provenance.

\begin{center}\small
\begin{tabularx}{\textwidth}{@{}lX@{}}
\toprule
Physical lines & Audited component\\
\midrule
665--770 & Original moment, fixed-character extraction, Mellin and Dirichlet deduction\\
1000--1248 & Initial Poisson reduction and original-family normalization\\
1261--1629 & Canonical estimate, cube comparison, transfer interface and finite induction\\
1649--2215 & Theta realization, local reflection, quadratic sieve and completed mean square\\
2217--2683 & Full two-Poisson transfer proof\\
2905--3478 & Fixed-ray theta construction, cusp multipliers, Mellin normalization and uniformity\\
\bottomrule
\end{tabularx}
\end{center}
The companion notes retain the separate September 30 correspondence.
That comparison does not assert that the two source texts or all
later claims in them are identical.

\section{Arithmetic family and the initial reduction}\label{sec:initial}

Write \(\OO=\mathbb Z[\omega]\), where
\(\omega=e^{2\pi i/3}\), and \(\lambda=1+2\omega=i\sqrt3\).
For an element \(a\), \(\N(a)=|a|^2\); for an ideal it is the
absolute ideal norm. The field has class number one. Ideals prime
to \(3\) are represented by their unique primary generators
\(a\equiv1\pmod3\). In row sums, elements of \(\OO\) are counted,
including their six unit associates; in starred column sums,
squarefree ideals prime to \(S\) are counted once.

For primary squarefree \(n\) prime to \(S\), let
\(\chi_n(x)=(x/n)_6\). Every local power is extended by zero
on nonunits:
\[
 \chi_p^j(x)=0\quad(p\mid x),\qquad
 \chi_p^0(x)=\ind_{p\nmid x}.
\]
The latter identity is essential: a zero local exponent does not
remove its original zero extension. When a nonsquarefree primary
ideal denominator occurs later, \(\chi_a\) is the product of its
local powers over the positive prime valuations, with each local
zero extension retained. The unit denominator has character one.
 Let \(\mu\) be the ideal
M\"obius function and define
\begin{equation}\label{eq:A}
 A_u(D)=\sum_{(n,S)=1}\mu(n)\nu(n)\chi_n(u)
                      W(\N(n)/D).
\end{equation}
Because \(\mu\) vanishes on nonsquarefree ideals, these are the
source's squarefree columns. The weight is complex-valued, smooth,
and supported in a fixed compact interval in \((0,\infty)\).

\subsection{The special Gauss-sum coefficients}

Set
\[
 e(z)=\exp(4\pi i\,\operatorname{Im}z/\sqrt3),\qquad
 \alpha(n)=n/|n|,\qquad
 \gamma_j(n)=\N(n)^{-1/2}\sum_{x\bmod n}\chi_n(x)^j e(x/n).
\]
For a fixed ray character \(\xi\), the returned coefficient is
\begin{equation}\label{eq:a}
 a_\xi(n)=\overline{\alpha(n)}\gamma_2(n)\xi(n).
\end{equation}
It has modulus one on the squarefree column domain. The finite
ray group is chosen before \(D\) and the moving primes; it accounts
for \(\nu\), reciprocity signs and supplementary Gauss phases.

\begin{assumption}[Arithmetic conversion interface]\label{in:arith}
The fixed ray group supports a unit-modulus function \(G\) and a
symmetric sign bicharacter \(\mathcal R\), with
\begin{align}
 \chi_b(a)&=\mathcal R(a,b)\chi_a(b),&
 G(ab)&=G(a)G(b)\mathcal R(a,b),\label{eq:recip-audit}\\
 \gamma_2(n)^3&=\mu(n)\alpha(n),&
 \gamma_1(n)\gamma_2(n)&=\mu(n)\alpha(n)G(n),\label{eq:gauss-audit}\\
 G(n)&=\overline{\chi_n(4)}\gamma_3(n),&
 \gamma_1(n)\gamma_{-1}(n)&=\chi_n(-1).
\end{align}
Here the residue-symbol reciprocity identities require the displayed
coprimalities, and the Gauss identities require squarefree \(n\)
prime to \(S\). The multiplicative identity for \(G\) is on the
fixed ray group. Consequently, for coprime squarefree \(a,b\),
\begin{align}
 a_\xi(ab)&=a_\xi(a)a_\xi(b)\chi_b(a)^4,\label{eq:crt-audit}\\
 \mu(n)\gamma_{-1}(n)
 &=\chi_n(-1)G(n)^{-1}a_1(n).\label{eq:conversion}
\end{align}
\end{assumption}

These are the source's finite Gauss--Jacobi and reciprocity
identities, audited with the inert-prime sign and the normalization
of \(e\) retained. The completed family below uses precisely
\eqref{eq:a}; no estimate for arbitrary sextic coefficients is inferred.

\begin{definition}[Dual energy]\label{def:energy}
For \(F\ge1\) and \(X,\mathcal H>0\), put
\begin{equation}\label{eq:energy-audit}
 \E(\mathcal H,X,F;\xi,W)=\frac1{XF}
 \sfstar_{F\le\N(f)<2F}\sum_{0<\N(k)\le\mathcal H}
 \left|\sfstar_n a_\xi(n)\chi_n(k)\chi_n(f)^4W(\N(n)/X)\right|^2.
\end{equation}
Write \(\E_r\) for the same energy with \((n,r)=1\).
\end{definition}

\begin{lemma}[Removal of column exclusions]\label{lem:exclusions}
Let \(r_0\) be the squarefree product of primes of \(r\) outside
\(S\). Then
\begin{equation}\label{eq:exclusions-audit}
 \E_r(\mathcal H,X,F;\xi,W)
 \le\tau(r_0)\sum_{d\mid r_0}
       \E(\mathcal H,X/\N(d),F\N(d);\xi,W).
\end{equation}
Every shift preserves \(XF\) and \(\mathcal H\).
\end{lemma}
\begin{proof}
Insert \(\sum_{d\mid(n,r_0)}\mu(d)\) in the column sum and put
\(n=dm\). Identity~\eqref{eq:crt-audit} supplies the factor
\(\chi_m(d)^4\), which enforces \((m,d)=1\); the twist with
\(f\) similarly makes terms with \((d,f)>1\) zero.
The remaining column has auxiliary twist \(df\) and scale
\(X/\N(d)\), times a coefficient of modulus at most one.
Cauchy--Schwarz over \(d\), followed by enlargement of the injective
squarefree image \(f\mapsto df\), gives \eqref{eq:exclusions-audit}.
The denominator remains
\((X/\N(d))(F\N(d))=XF\). No vanishing character value is divided out.
\end{proof}

\subsection{Excluded Poisson and normalization}

For a primitive finite character \(\chi\) of modulus \(\mathfrak m\),
the normalized excluded Poisson formula used in the source is
\begin{equation}\label{eq:poisson-audit}
\begin{split}
 \sum_k\chi(k)\ind_{(k,r)=1}\Phi(\N(k)/\mathcal H)
 ={}&\frac{\mathcal H\gamma(\chi)}{\sqrt{\N(\mathfrak m)}}
 \sum_{d\mid\operatorname{rad}r}\frac{\mu(d)\chi(d)}{\N(d)}\\
 &\hspace{4mm}\times\sum_h\overline{\chi(h)}
  \widehat\Phi\!\left(\frac{\mathcal H\N(h)}
                           {\N(d)\N(\mathfrak m)}\right).
\end{split}
\end{equation}
The lattice measure is normalized to give \(\OO\) covolume one.
For the principal primitive character, the modulus is \(1\) and
the zero frequency is
\(\mathcal H\widehat\Phi(0)\prod_{p\mid r}(1-\N(p)^{-1})\).
For nonprincipal primitive characters the zero frequency vanishes.
This analytic identity is a classical Poisson input, with the
primitive finite Fourier identity specifying its normalization.

Choose nonnegative radial \(\Phi\ge1\) on \([0,1]\) with compact
Fourier support. The initial smoothed moment is
\begin{equation}\label{eq:MD}
 \mathcal M_D=\frac1D\sum_{u\in\OO}\Phi(\N(u)/H)|A_u(D)|^2,
 \qquad H=D^{1+\vartheta}.
\end{equation}
In the expanded columns write \(n_j=gz_j\), with
\((z_1,z_2)=(g,z_1z_2)=1\). The common column factor produces
\(\ind_{(u,g)=1}\), which belongs in \eqref{eq:poisson-audit}.
The zero frequency occurs only for \(z_1=z_2=1\), and gives
\[
 Z=\frac HD\widehat\Phi(0)\sfstar_g |W(\N(g)/D)|^2
       \prod_{p\mid g}(1-\N(p)^{-1})
 \ll_I H\norm{W}_\infty^2.
\]
Nonzero frequencies with \(z_1=z_2=1\) remain in the off-zero
expression; deleting them would not be a diagonal subtraction.

After inserting the remaining coprimality M\"obius sum, the exact
change of variables is
\[
 b=g/e,\quad f=ev,\quad k=eh,\qquad
 e=(f,k),\quad v=f/e,\quad g=be,\quad h=k/e.
\]
It yields \(gv=bf\), \(he^5v^4=kf^4\),
\(\mu(e)\mu(v)=\mu(f)\), and \(\N(g)/\N(e)=\N(b)\).
The common-factor masks are part of the inverse map. The paired
Gauss conversion and fixed ray expansion return \(a_\xi\) and its
conjugate, while the column exclusion left over is \((n,b)=1\).

\begin{assumption}[Initial analytic reduction]\label{in:initial}
The source reduction, with these masks and the full diagonal,
has the following interface. For fixed \(C\ge1\), dyadic \(B,F\ge1\)
give
\begin{equation}\label{eq:initial-ranges}
 X=D/(BF),\qquad
 0<\mathcal H\le CD^{1-\vartheta}/B^2,\qquad
 XF=D/B.
\end{equation}
A bound
\(\E(\mathcal H,X,F;\xi,U)
 \ll D^\epsilon XF\norm{U}_{C^J}^2\)
for every common test implies \eqref{eq:main-moment}.
The row-dependent two-column kernel is controlled by its
\(C^{2J+4}\) norm, bounded by
\(C_{I,J}\norm{W}_{C^{2J+4}(I)}^2\).
\end{assumption}

The source's smooth-weight lemma is used only after the common-test
bound is available. The final recovery, including short divisor shifts
and the outside factor \(D\), is proved explicitly in
Section~\ref{sec:moment-final}.

\section{Theta reflection and the completed mean square}\label{sec:completed}

\subsection{Completed family and reflected interface}

For a completely multiplicative ideal function \(\Psi\), zero at
primes in \(S\), let \(V_*(y)=\sqrt y\,W(y)\) and define
\begin{equation}\label{eq:T-audit}
 T(X;\Psi)=\sfstar_n\sum_{\substack{b\equiv1\ (3)\\(b,S)=1}}
 \frac{\overline{\alpha(n)}\gamma_2(n)\Psi(n)
        \overline{\alpha(b)}^{\,3}\Psi(b)^3}
      {\sqrt{\N(n)}\N(b)}
 V_*(\N(n)\N(b)^3/X).
\end{equation}
Only \(n\) is squarefree. For the energy family,
\(\Psi_{k,f}(n)=\xi(n)\chi_n(k)\chi_n(f)^4\) on \((n,S)=1\)
and zero otherwise; write \(T(X;k,f)\). The term \(b=1\)
is \(X^{-1/2}\) times the uncompleted column.

The angular Mellin calculation gives the transform
\begin{equation}\label{eq:sharp}
 V_*^\sharp(x)=\frac1{2\pi i}\int_{(0)}
 \widehat V_*(-t)
 \frac{\Gamma(7/6+t)\Gamma(5/6+t)}
      {\Gamma(7/6-t)\Gamma(5/6-t)}
 \left(\frac{(2\pi)^4x}{27}\right)^{-t}\dd t.
\end{equation}
The reflected kernel uses \(\N(\ell)X/\N(c)^2\).
Thus its denominator is \(\N(c)^2=|c|^4\), not \(\N(c)\).
The scalar is \(-i/81\) times the bounded Fourier and arithmetic
phases. These normalizations were independently derived within
the audit from the differentiated theta Mellin integral.

\begin{assumption}[Theta reflection and moving-row uniformity]\label{in:theta}
For
\(\Psi=\Psi_0\prod_{p\in\mathcal P}\chi_p^{j_p}\),
\(0\le j_p\le5\), the completed sum is a finite sum of terms
\begin{equation}\label{eq:reflected}
 C\sum_{\ell\ne0}\frac{d(\ell)\alpha(\ell)}{\sqrt{\N(\ell)}}
 \psi(\lambda^4\ell)
 \prod_{p\in\mathcal A}B_{p,j_p}(\lambda^4\ell)
 V_*^\sharp\!\left(\frac{\N(\ell)X}{\N(c)^2}\right),
\end{equation}
where \(\{p:j_p\ne0\}\subseteq\mathcal A\subseteq\mathcal P\),
\(c=c_0\prod_{p\in\mathcal A}p\), and \(C\) is bounded independently
of moving prime norms. The local factors are
\begin{equation}\label{eq:B}
 B_{p,j}(x)=
 \begin{cases}
  \chi_p(x)^{-j-2},&j\notin\{0,4\},\\
  q^{-1/2}(-1+q\ind_{p\mid x}),&j=4,\\
  q^{-1/2}\chi_p(x)^{-2},&j=0,
 \end{cases}\qquad q=\N(p).
\end{equation}
The cusp coefficient \(d\), additive character \(\psi\), and
fixed denominator \(c_0\) lie in a fixed finite family. For residual
squarefree moving row \(k_0\), every prime of \(k_0\) is active with
exponent one. After partitioning \(k_0\) into fixed ray classes,
\((d,\psi,c_0)\) is independent of the moving row within a sector.
The moving denominator still contains \(k_0\).

The coefficient support and bounds are
\begin{equation}\label{eq:theta-bounds-audit}
 \ell=u\lambda^mnb^3,\quad m\ge-4,\quad n\text{ squarefree},
 \qquad |d(\ell)|\le27\,3^{m/6}\sqrt{\N(b)}.
\end{equation}
Here \(u\) is a unit and \(n,b\) are primary and prime to
\(3\). The original \(S\)-exclusion is not reimposed on dual columns,
and there is no condition \((n,b)=1\). Finally,
\begin{equation}\label{eq:decay}
 |(x\partial_x)^jV_*^\sharp(x)|
 \ll_{A,j,I}\norm{W}_{C^J(I)}
              \min(x^{1/4},x^{-A}),\qquad
 J=\lceil4A\rceil+j+4
\end{equation}
is a sufficient finite-order bound for every fixed \(A>0,j\ge0\).
\end{assumption}

The cusp formulas were compared with Dunn--Radziwi{\l}{\l}
\cite{DR}. Their coefficients are conjugated and their frequencies
reversed in \eqref{eq:reflected}; ignoring this distinction changes
the two nonidentity cusps. The original Patterson proof \cite{Pat}
was not independently read. The source's reduced-denominator
construction, which permits nonunit fixed Fourier frequencies,
was audited separately; the coprime-frequency statements in
\cite{DR} alone do not supply that extension.

Two analytic clarifications matter. Cusp representatives and CRT
solutions can be moved within their allowed congruences to avoid
the finitely many zero matrix entries used in residue symbols.
For a fixed arithmetic twist, differentiating the theta function
removes the constant mode; exponential decay at both height
endpoints makes its Mellin transform entire. The direct and dual
series give polynomial bounds on strip edges. A crude interior
exponential bound, together with
\(\exp(\eta s^2)(B+s)^{-M}\) on expanding rectangles, gives
polynomial strip growth by the maximum principle and then
\(\eta\downarrow0\). Kernel contours are shifted individually before
their rapidly decaying values are summed; an unweighted dual
series is not interchanged with the zero-line integral.

The gamma quotient in \eqref{eq:sharp} has no poles in
\(-1/4\le\Re t\le A\) and has polynomial growth
\(O_A((1+|\operatorname{Im}t|)^{4\Re t})\).
Compact support and finitely many integrations by parts therefore
give \eqref{eq:decay}, uniformly in the moving primes. Constants
used to establish the exact fixed-twist identity may depend on
that twist; they are not subsequently claimed as uniform estimates.

\subsection{Exceptional primes and their exact cost}

For a squarefree row write \(s=tk_0\), with \(t\) supported on
the primes of \(gS\), and put \(\mathcal H_0=\mathcal H/\N(t)\).
The original sextic twist is \(k=u_0tk_0\) with auxiliary \(g\),
not necessarily squarefree. Away from \(S\) its local exponent is
\(j_p=v_p(k)+4v_p(g)\pmod6\).
If \(p\mid k_0\), then \(j_p=1\) and
\(B_{p,1}(nb^3)=\chi_p(nb)^3\), including its zeros.
If \(p\mid t\) with \(p\notin S\), the exponent
\(1+4v_p(g)\) is odd; a shared prime of \(t\) and \(g\)
must not be charged as a second exceptional prime. Primes in
\(S\) belong to the fixed denominator and finite-sector data,
rather than this sextic local table.

At a prime of norm \(q\) with \(j=4\), the exact disjoint split is
\begin{equation}\label{eq:j4}
 q^{-1/2}(-1+q\ind_{p\mid nb})
 =-q^{-1/2}+q^{1/2}\ind_{p\mid n}
                    +q^{1/2}\ind_{p\nmid n,\ p\mid b}.
\end{equation}
A collision \(p\mid n,b\) belongs to the second term once.
Reindexing \(n=pn'\) requires \(p\nmid n'\); reindexing \(b=pb'\)
requires \(p\nmid n\), with no restriction on further powers in
\(b'\). The coefficient at all other primes is evaluated at the
original substituted pair. This argument uses its magnitude bound,
not multiplicativity of the cusp coefficient.

Let a prepared coefficient have bound \(27a\) and auxiliary length
\(Y\). The complete local cost is:
\begin{center}\small
\begin{tabular}{@{}lccc@{}}
\toprule
Branch & New \(a^2\) & New \(Y\) & Cost in \(a^2Y\)\\
\midrule
Active \(j=1,2,3,5\) & \(a^2\) & \(q^2Y\) & \(q^2\)\\
Inactive \(j=0\) & \(\le a^2\) & \(Y\) & \(\le1\)\\
Active \(j=0\) & \(a^2/q\) & \(q^2Y\) & \(q\)\\
\(j=4\), constant & \(a^2/q\) & \(q^2Y\) & \(q\)\\
\(j=4\), squarefree column & \(a^2\) & \(qY\) & \(q\)\\
\(j=4\), cube column & \(a^2/q\) & \(Y/q\) & \(q^{-2}\)\\
\bottomrule
\end{tabular}
\end{center}
The inactive zero-exponent branch has scalar \(1-q^{-1}\) and
no transformed column mask at \(p\). Restoring the original mask
after reflection would change the exact identity.

Starting at \(Y=\N(c_0)^2\mathcal H_0^2/X\), the table gives
\begin{equation}\label{eq:cost}
 a^2\le1,\qquad
 a^2Y\ll_{\xi,S}\mathcal H_0^2\N(t)^2\N(g)/X.
\end{equation}
Indeed, primes of \(t\) cost \(q^2\). For a prime of \(g\) outside
\(t\), valuation \(r\equiv0,1,2\pmod3\) gives respectively
\(j=0,4,2\), with cost at most \(q,q,q^2\), bounded by \(q^r\).
The number of branches is divisor bounded in the polynomial norm
range. Fixed ray partitioning, zero masks, fixed substitutions and
bounded exterior row phases leave each prepared coefficient
independent of \(k_0\).

\subsection{The primitive quadratic family}

For squarefree primary \(k\) prime to \(2,3\), define
\[
 e_k=\begin{cases}0,&\N(k)\equiv1\pmod4,\\1,&\N(k)\equiv3\pmod4,\end{cases}
 \qquad
 \Xi_k(x)=(x/k)_2\kappa_\lambda(x)^{e_k},
\]
with \(\kappa_\lambda\) the nontrivial quadratic character modulo
\(\lambda\). The supplement is absent when \(e_k=0\).
Since \(\omega\) is a square unit and
\((-1/k)_2=(-1)^{e_k}\), the supplement makes \(\Xi_k\)
trivial on the units. Its primitive conductor is
\((k)\lambda^{e_k}\), its infinite type is trivial, and its order
divides two. There is no added conductor at \(2\) in this realization.
For a primary column \(n\), \(\kappa_\lambda(n)=1\), so
\(\Xi_k((n))=\chi_k(n)^3\).

Partition into the fixed reciprocity classes, including the parity
of \(\N(k)\) modulo \(4\). For coprime \(k_1,k_2\) in the same
class, the supplements cancel, and the primitive inducing product
has conductor \(k_1k_2\). This is the inducing-product convention
in Goldmakher--Louvel \cite{GL}, not a claim that every
zero-extended product at \(\lambda\) has that conductor.
Their quadratic-family theorem, together with fixed-class
reciprocity or duality, gives
\begin{equation}\label{eq:Q-audit}
 \sfstar_{\N(k)\le\mathcal H_0}
 \left|\sfstar_{\N(n)\asymp U}\beta(n)\chi_k(n)^3\right|^2
 \ll_\sigma(\mathcal H_0U)^\sigma
                    (\mathcal H_0+U)\sum_n|\beta(n)|^2.
\end{equation}
Dual columns meeting fixed \(S\)-primes are handled by writing
\(n=s_0n_0\), with \(s_0\) a squarefree \(S\)-supported ideal.
There are finitely many choices; its character value is a bounded
row scalar. This permits, in particular, the cubic theta coefficient
at the norm-\(4\) prime without extending the original sextic
denominator notation there.

\subsection{Common Mellin majorant and all infinite sums}

In a fixed ray sector, unit \(u\), and ramified index \(m\ge-4\),
the prepared sum has the form
\[
 \mathscr S_m(k_0)=
 \sum_{n,b}\frac{A_m(n,b)\chi_{k_0}(nb)^3}
                   {3^{m/3}\sqrt{\N(n)}\N(b)}
 V_*^\sharp\!\left(
  \frac{3^m\N(n)\N(b)^3\mathcal H_0^2}{Y\N(k_0)^2}\right),
 \quad |A_m(n,b)|\le27a .
\]
For dyadic \(\N(n)\asymp U,\N(b)\asymp B\), let
\(z_*=3^mUB^3/Y\). Mellin inversion in the column variable
has the common pointwise bound
\[
 |c_{b,k_0}(\tau)|\ll
 \norm{W}_{C^J}(1+z_*)^{-A}(1+|\tau|)^{-2}.
\]
The lower bound for the kernel argument uses
\(\N(k_0)\le\mathcal H_0\). A separate bounded integral for each
row would not justify exchanging the moving row sum with Mellin;
the displayed common bound does.

For the sieve coefficients
\(\beta(n)=A_m(n,b)\N(n)^{-1/2}(\N(n)/U)^{i\tau}\),
ideal counting gives \(\sum_n|\beta(n)|^2\ll a^2\).
Weighted Cauchy in Mellin keeps two factors of kernel decay.
Cauchy in the cube index uses
\(\sum_{\N(b)\asymp B}\N(b)^{-1}\ll1\), leaving no extra \(B\).
Together with \eqref{eq:Q-audit}, this gives
\begin{equation}\label{eq:block-completed}
 \sum_{k_0}|\mathscr S_{m;U,B}(k_0)|^2
 \ll a^2 3^{-2m/3}\norm{W}_{C^J}^2
  (\mathcal H_0U)^\sigma(\mathcal H_0+U)
  (1+3^mUB^3/Y)^{-2A}.
\end{equation}

The needed dyadic sums, for every \(Z>0\), are
\[
 \sum_{U,B\ge1}(1+UB^3/Z)^{-2}\ll\log^2(2+Z),\qquad
 \sum_{U,B\ge1}\sqrt U(1+UB^3/Z)^{-2}\ll\sqrt Z .
\]
For \(Z\le1\), the left sides are \(O(Z^2)\), by geometric sums
of \(U^{-2}B^{-6}\) and \(U^{-3/2}B^{-6}\). For \(Z\ge1\),
the region \(UB^3\le Z\) contributes the stated orders and the
complement has geometric decay. Thus no restriction \(Y\ge1\)
is introduced. For \(z_*=3^mUB^3/Y\), one has
\(U\le\max(1,Y/3^m)(1+z_*)\). Since \(m\ge-4\) and \(Y\)
has a fixed polynomial bound, the small sieve factor \(U^\sigma\)
is absorbed by a reserved power of \(D\) and a small reduction in
the decay exponent; no unbounded \(U\) is simply declared polynomial.
Minkowski applied to \eqref{eq:block-completed},
with reserved small losses and \(A\ge3\), gives for each prepared
term the bound \(D^\epsilon(\mathcal H_0a^2+a^2Y)\)
after summing \(U,B\) and \(m\ge-4\).
The row part has geometric factor \(3^{-m/3}\), while the column
part has \(3^{-5m/6}\sqrt Y\), with only finitely many negative \(m\).
Polynomial bounds for \(Y\), branch counts and logarithms allow
all the small losses to be chosen before summation.
For this separation and decay step, \(J\ge18\) is a sufficient
order; later compositions may require a larger finite order.

Using \eqref{eq:cost} and summing the divisor-bounded \(t\)-sectors
therefore gives the squarefree-row bound
\begin{equation}\label{eq:sf-completed}
 \sum_{\substack{0<\N(s)\le\mathcal H\\s\ {\rm squarefree}}}
 |T(X;u_0s,g)|^2
 \ll D^\epsilon\norm{W}_{C^J}^2
             \left(\mathcal H+\frac{\mathcal H^2\N(g)}X\right)
\end{equation}
in any fixed polynomial parameter range, for arbitrary nonzero
\(g\), including its repeated factors and primes in \(S\).

\begin{proposition}[General rows in the completed family]\label{prop:completed}
Under the reflected and quadratic interfaces above, for
\(1\le\mathcal H,X,\N(f)\le D^{C_0}\) and squarefree primary
\((f,S)=1\),
\begin{equation}\label{eq:completed-bound}
 \sum_{0<\N(k)\ll\mathcal H}|T(X;k,f)|^2
 \ll_{I,\xi,S,\epsilon,C_0}D^\epsilon\norm{W}_{C^J(I)}^2
                 \left(\mathcal H+\frac{\mathcal H^2\N(f)}X\right).
\end{equation}
The order \(J=J(\epsilon,C_0)\) is finite and independent of \(I\).
\end{proposition}
\begin{proof}
Write uniquely \(k=u_0sv^2\), with \(s\) squarefree and one
multiplicatively chosen generator per ideal for \(s,v\).
The supports of \(s,v\) may overlap. Complete multiplicativity,
including zero values, and \(8\equiv2\pmod6\) give
\[
 T(X;u_0sv^2,f)=T(X;u_0s,fv^2).
\]
For each \(v\), let \(\mathcal H_v=\mathcal H/\N(v)^2\) and
\(g=fv^2\). A nonempty range has \(\mathcal H_v\ge1\) and
\(\N(g)\le\mathcal H\N(f)\le D^{2C_0}\), so
\eqref{eq:sf-completed} applies with range \(2C_0\). Its term is
exactly
\[
 \mathcal H_v+\mathcal H_v^2\N(g)/X
 =\frac1{\N(v)^2}
        \left(\mathcal H+\frac{\mathcal H^2\N(f)}X\right).
\]
The ideal sum \(\sum_v\N(v)^{-2}=\zeta_K(2)\) converges.
The six units cost a fixed factor. This proves the proposition
without an extra count of \(v\) or a restriction \((s,v)=1\).
\end{proof}


% Prepared for Edward Baker, 12 October 2026, with substantial LLM assistance.
% GPT-6 (Codex), inherited configuration; exact variant/effort unavailable.
% Conditional reconstruction from the pinned source and three transfer notes.
% Fragment only: the parent manuscript supplies packages, macros and bibliography.

\section{Two Poisson transformations and finite descent}
\label{sec:consolidated-transfer}

We retain the field, fixed excluded set $S$, ray characters $\xi$, and
coefficients $a_\xi(n)$ of the preceding sections. Ideal indices are
represented by primary generators; a star means squarefree and prime to
$S$. Rows are actual elements of $\OO$, with their unit associates retained.
All residue characters have their zero extensions. The present section
reconstructs the return and descent argument of \cite{Source}, physical
lines 1335--1629 and 2217--2683. Its arithmetic Poisson and completed-sum
inputs remain identified interfaces, rather than independently proved
theta or large-sieve theorems.

\begin{assumption}[Interfaces used in the return argument]
\label{ass:transfer-interfaces}
We use the source's first positive-form and second Poisson identities,
with their full enlarged outer domains, common kernels, zero modes
and arbitrary positive intermediate scales, as described below.
Their analytic identities are source interfaces; the signed
regrouping and admissible-scale algebra following them are derived
in this section.

The coefficients have modulus one on the starred domain and obey
$a_\xi(cn)=a_\xi(c)a_\xi(n)\chi_n(c)^4$ on coprime squarefree indices.
We use the normalized excluded Poisson formula \eqref{eq:poisson-audit}
for every positive row scale, including its principal modulus-one case.
The Fourier and primitive Gauss normalization is retained. The paired Gauss
conversions on coprime squarefree columns are
\[
\begin{aligned}
 a_\xi(u_1)\overline{a_\xi(u_2)}\gamma(\chi_{u_1}\overline{\chi_{u_2}})
 &=\mu(u_1)\mu(u_2)\sum_{\xi_1}\ell_{\xi_1}
       \xi_1(u_1)\overline{\xi_1(u_2)},\\
 \mu(z_1)\mu(z_2)\xi_1(z_1)\overline{\xi_1(z_2)}
 \gamma(\overline{\chi_{z_1}}\chi_{z_2})
 &=\sum_{\xi'}c_{\xi'}a_{\xi'}(z_1)\overline{a_{\xi'}(z_2)}.
\end{aligned}
\]
The character sets and coefficients are finite and fixed once $\xi$ or
$\xi_1$ is fixed; they do not depend on the moving ideal or row indices.
These are the source's CRT, reciprocity and Gauss conversion interfaces.

We also use the following common-test smooth interface. For a finite row
set with $w_r\ge0$, put $S_{j,r}(V)=\sum_n a_{j,r}(n)V(x_{j,r,n})$
with finite column sums, and suppose both families satisfy
$\sum_r w_r|S_{j,r}(V)|^2\le M_j\norm{V}_{C^m(J_*)}^2$ for every
common test $V$, where $J\subset\operatorname{int}J_*$. If
$|b_r|\le w_r$ and $K_r$ is supported in $J^2$, then
\[
 \left|\sum_r b_r\sum_{n_1,n_2}
 a_{1,r}(n_1)\overline{a_{2,r}(n_2)}K_r(x_{1,r,n_1},x_{2,r,n_2})\right|
 \ll\sqrt{M_1M_2}\sup_r\norm{K_r}_{C^{2m+4}(J^2)}.
\]
This is the row-dependent Mellin separation interface of the source,
lines 3570--3650: its Fourier coefficients have one common pointwise
majorant before rows are summed.

Finally, with $\Psi_k(n)=\xi(n)\chi_n(k)\chi_n(f)^4$ and its fixed
$S$-zero extension, the completed family of the preceding section is
\[
 T(X;k,f)=\frac1{\sqrt X}\sum_n^*\sum_{(b,S)=1}
 \sqrt{\N(b)}\,\overline{\alpha(b)}^{\,3}\Psi_k(b)^3
 a_\xi(n)\chi_n(k)\chi_n(f)^4W(\N(n)\N(b)^3/X),
\]
and satisfies
\[
 \sum_{0<\N(k)\le H}|T(X;k,f)|^2
 \ll D^\eta\norm{W}_{C^{J_R}(I)}^2
       \left(H+\frac{H^2\N(f)}X\right)
\]
for $1\le H,X,\N(f)\le D^C$, every $\eta>0$, and
$J_R=J_R(\eta,C)$ independent of $I$. The theta transformation and
quadratic sieve underlying this estimate are the preceding section's
deep inputs. No arbitrary sextic coefficient family is substituted.
\end{assumption}

Write $\E_{(r)}$ for the normalized energy with the extra column condition
$(n,r)=1$. Its normalization is $1/(XF)$. Inclusion--exclusion and the CRT
law give, for squarefree $r$ prime to $S$,
\begin{equation}\label{eq:transfer-exclusions}
 \E_{(r)}(\mathcal H,X,F;\xi,V)
 \le\tau_{\rm div}(r)\sum_{j\mid r}
 \E(\mathcal H,X/\N(j),F\N(j);\xi,V).
\end{equation}
As in Lemma~\ref{lem:exclusions}, $n=jm$ changes $f$ to $jf$ without
changing the test; Cauchy--Schwarz in $j$ and enlargement of that injective
image prove the inequality. The product $(X/\N(j))(F\N(j))=XF$ is unchanged.

\subsection{The two-Poisson return estimate}
Fix a nonnegative radial Schwartz $\Phi$, with $\Phi\ge1$ on $[0,1]$ and
$\operatorname{supp}\widehat\Phi\subset[0,C_\Phi]$. For $I=[u,v]$, set
$I_*=[u/2,2v]$ and $I'=[u/16,4v]$. Define
\[
 \mathcal A(W)=\frac1{LF}\sum_{F\le\N(f)<2F}^*
 \sum_{k\in\OO}\Phi(\N(k)/\mathcal H)
 \left|\sum_n^*a_\xi(n)\chi_n(k)\chi_n(f)^4W(\N(n)/L)\right|^2.
\]
In particular $\E(\mathcal H,L,F;\xi,W)\le\mathcal A(W)$.

\begin{proposition}[Return to the same arithmetic family]
\label{prop:consolidated-transfer}
Let $D\ge2$ and $C_0\ge1$. Suppose $1\le\mathcal H,L,F,\Sigma\le D^{C_0}$ and
$\max\{\mathcal H,LF\}\le\Sigma$, where the target $\Sigma$ is independent
of $LF$. For every integer $m\ge0$ and $\epsilon>0$,
\[
 \mathcal A(W)\ll D^\epsilon\Sigma
       (1+S_m)\norm{W}_{C^{4m+12}(I)}^2.
\]
Here $S_m$ is the supremum of
$\E(\mathcal H',X',F';\xi',V)/\Sigma'$, with $\Sigma'=X'F'$,
$\mathcal H',X',F'\ge1$, $V$ supported in $I'$ and
$\norm{V}_{C^m(I')}\le1$, over precisely the restrictions
\begin{equation}\label{eq:consolidated-child}
 \mathcal H'\le\frac{\mathcal HL}{\Sigma F},\qquad
 \frac{\mathcal H'}{\Sigma'}\le\frac{\mathcal H}{\Sigma},\qquad
 \Sigma'\le L.
\end{equation}
An empty supremum is zero. Constants may depend on $I,m,C_0,\epsilon$
and the fixed arithmetic data, but not on moving scales or indices.
The returned $\xi'$ ranges over the fixed finite set of ray characters.
\end{proposition}

\begin{proof}
Choose $c_I\ge\max(1,4C_\Phi v^2)$. For disjoint squarefree $C,t$ and
$d\mid C$, put
\[
 \ell=\frac L{\N(C)\N(t)},\quad
 w_{C,d}=\frac{\mathcal H\N(C)}{\N(d)L^2F},\quad
 Y_{C,d}=\frac{c_I\Sigma LF\N(d)}{\mathcal H\N(C)^2},
\]
and retain the full outer domain $\N(Ct)\le2vL$. Define
\[
\begin{aligned}
 p_y(n)&=\mu(n)\xi_1(n)\mathbf1_{(n,t)=1}
       \overline{\chi_n(y)}\chi_n(C)^4\chi_n(d),\\
 \mathcal Q_{\xi_1}(U)&=\sum_{C,t}\sum_{d\mid C}w_{C,d}
       \sum_{y\in\OO}\Phi(\N(y)/Y_{C,d})
       \left|\sum_n^*p_y(n)U(\N(n)/\ell)\right|^2.
\end{aligned}
\]
The outer sums are finite and the $y$-sum is absolutely convergent.
No first-stage Fourier-support restriction is imposed on this outer domain.

For the first Poisson application extract $C=(n_1,n_2)$, write $n_i=Cu_i$,
and use the first paired conversion in Assumption~\ref{ass:transfer-interfaces}.
Insert $\sum_{t\mid u_1,u_2}\mu(t)$ and put $u_i=tx_i$.
The common factors give exactly $(f,Ct)=(h,t)=1$; individual column zeros
remain in $p_{hf^2}$. The nonzero contribution is a finite signed sum,
with exterior weight $w_{C,d}$ and coupled kernel
\[
 (z_1z_2)^{-1/2}W(z_1)\overline{W(z_2)}
 \widehat\Phi\!\left(\frac{\mathcal H\N(h)\N(C)^2}
 {L^2\N(d)z_1z_2}\right),\qquad z_i=\N(x_i)/\ell.
\]
This is the expanded identity in \cite{Source}, lines 2293--2327.
Its $C^q$ norm is $\ll\norm{W}_{C^q(I)}^2$ uniformly on the finite
frequency range. The zero Fourier mode requires $u_1=u_2=1$ and costs
$O_I(\mathcal H\norm{W}_\infty^2)$; all principal nonzero frequencies remain.
The map $(f,h)\mapsto y=hf^2$ has divisor-bounded multiplicity and
$0<\N(y)\le Y_{C,d}$, since $\Sigma\ge LF$. Its image norms are polynomial
in $D$. If $Y_{C,d}<1$ the nonzero image is empty. Positivity and the smooth
interface therefore imply that a common bound
$\mathcal Q_{\xi_1}(U)\le M\norm{U}_{C^j(I_*)}^2$ yields
\begin{equation}\label{eq:transfer-first-bound}
 \mathcal A(W)\ll D^{\epsilon_0}(\Sigma+M)
                  \norm{W}_{C^{2j+4}(I)}^2.
\end{equation}

Apply Poisson again to $\sum_y\Phi(\N(y)/Y)|P(y)|^2$, initially for any
$\ell,Y>0$. Extract the common column $g$, take $e\mid g$ from its row
exclusion, and insert the second column divisor $w$. The second paired
conversion returns the coefficient
\[
 a^\sharp(n)=a_{\xi'}(n)\mathbf1_{(n,tg/e)=1}
                      \chi_n(deh)\chi_n(Cew)^4.
\]
The ranges are $(g,Ct)=(w,gCth)=1$,
$\N(gw)\le2v\ell$, and
$0<\N(h)\le4C_\Phi v^2\ell^2\N(e)/(Y\N(g)^2)$.
With $U_0(x)=x^{-1/2}U(x)$, the exterior coefficient and kernel are
$Y\mu(e)\mu(w)c_{\xi'}\N(g)/(\N(e)\ell)$ and
$U_0(x_1)\overline{U_0(x_2)}
\widehat\Phi(Y\N(h)\N(g)^2/(\N(e)\ell^2x_1x_2))$.
The individual column scale is $\ell/\N(gw)$. This exact finite output
is the second Poisson interface at source lines 2407--2494. Its residual
mask follows from zero extension, including
$\chi_z(dh)\overline{\chi_z(e)}\chi_z(C)^4
=\chi_z(deh)\chi_z(Ce)^4$; no inverse of a zero character value is used.
Its zero Fourier term is
\[
 Z_2=Y\widehat\Phi(0)\sum_{(g,Ct)=1}^*
 |U(\N(g)/\ell)|^2\prod_{p\mid g}(1-\N(p)^{-1})
 \ll Y\ell\norm{U}_\infty^2.
\]
This holds also for $Y<1$: if $2v\ell<1$ columns are empty, and otherwise
their number is $O_I(\ell)$. It does not bound the entire row sum by $O(Y)$.

Now substitute $Y=Y_{C,d}$ and set $r=tg/e$, $f'=Cew$, $k'=deh$.
Then $r,f'$ are disjoint squarefree indices, and exactly
\begin{equation}\label{eq:transfer-common-normalization}
 X'_0=\frac L{\N(r)\N(f')},\quad
 w_{C,d}\frac{Y_{C,d}\N(g)}{\N(e)\ell}
 =\frac{c_I\Sigma\N(r)}{L^2},\quad
 \frac{Y_{C,d}\N(h)\N(g)^2}{\N(e)\ell^2}
 =\frac{c_I\Sigma F\N(k')\N(r)^2}{\mathcal HL}.
\end{equation}
Thus all preimages have the same kernel and unsigned normalization.
For fixed $r,f',k'$, every preimage is obtained by choosing
$t\mid r$, a disjoint factorization $f'=Cew$, $e\mid k'$,
$(w,k')=1$, and $d\mid(C,k')$, then taking $g=e(r/t)$ and $h=k'/(de)$.
These choices satisfy every old coprimality condition. Since
$Ctgw=rf'$, all belong to the enlarged outer domain whenever the common
kernel is nonzero; the frequency conditions also become the same common bound.

The $t$ choices contribute $\tau_{\rm div}(r)$. For $p\mid f'$ put
$\delta_p=\mathbf1_{p\mid k'}$. Placement in $C,e,w$ contributes respectively
$(1+\delta_p),-\delta_p,-(1-\delta_p)$ to the signed sum of
$\mu(e)\mu(w)$, whose local total is $\delta_p$. Therefore
\[
 \sum_{\text{preimages}}\mu(e)\mu(w)
 =\tau_{\rm div}(r)\mathbf1_{f'\mid k'}.
\]
This includes arbitrary valuations in $k'$, its unit factor, intersections
with $r$, and unit ideal indices. Regrouping is performed before absolute
values. Only afterwards may positive mean squares be enlarged.

Nonzero evaluated terms satisfy $\N(rf')\le2vL$ and
$0<\N(k')\le\mathcal HL/(\Sigma F\N(r)^2)$.
In a dyadic block $R\le\N(r)<2R$, $F'\le\N(f')<2F'$, put
\[
 X'=\frac L{RF'},\qquad\Sigma'=\frac LR,\qquad
 \mathcal H'=\frac{\mathcal HL}{\Sigma FR^2}.
\]
Retain only $r$ with an actual surviving term in this block, before
enlarging $f',k'$. Its witness has
$F'\le\N(f')\le\N(k')\le\mathcal HL/(\Sigma F\N(r)^2)$.
Consequently $\mathcal H'\ge1$ and, for every $j\mid r$,
\[
 \frac{X'}{\N(j)}\ge
 \frac{\Sigma F\N(r)^2}{\mathcal HR\N(j)}
 \ge\frac{\Sigma F}{\mathcal H}\ge1.
\]
The exclusion-shifted child
$(\mathcal H',X'/\N(j),F'\N(j))$ therefore has all scales at least one,
preserves $\Sigma'$, and satisfies \eqref{eq:consolidated-child}.
Its individual scales are bounded by $D^{C_0}$. No additional $r$ is admitted.

The ratio $s=X'/X'_0$ lies in $[1,4)$, so the rescaled kernels are supported
in $[u/8,2v]^2\subset\operatorname{int}(I'^2)$ and have uniformly bounded
$C^{2m+4}$ norm $\ll\norm{U}_{C^{2m+4}(I_*)}^2$. The Fourier parameter
is bounded on nonzero support; there is no $Y^{-1}$ or divisor cost in this norm.
For each retained $r$, \eqref{eq:transfer-exclusions} and $S_m$ bound the
unnormalized common-test square sum by
$D^\delta(\Sigma')^2S_m\norm{V}_{C^m(I')}^2$.
The square of $\Sigma'$ includes undoing the energy normalization.
There are $O(R)$ retained $r$ and their exterior factor costs
$\Sigma R/L^2$ times a divisor loss. Smooth separation bounds this block by
\[
 D^{2\delta}\frac{\Sigma R}{L^2}R(\Sigma')^2S_m
 \norm{U}_{C^{2m+4}(I_*)}^2
 =D^{2\delta}\Sigma S_m\norm{U}_{C^{2m+4}(I_*)}^2.
\]
The total second zero modes use the exact identity
$w_{C,d}Y_{C,d}\ell=c_I\Sigma/(\N(C)^2\N(t))$ and cost
\[
 \ll\Sigma\norm{U}_\infty^2
 \sum_C\frac{\tau_{\rm div}(C)}{\N(C)^2}
 \sum_{\N(t)\le2vL}\frac1{\N(t)}
 \ll D^\delta\Sigma\norm{U}_\infty^2.
\]
The first series converges and the second is logarithmic. The number of
dyadic blocks is $O_I((\log(2D))^2)$; taking $\delta=\epsilon_0/4$ gives
$\mathcal Q_{\xi_1}(U)\ll D^{\epsilon_0}\Sigma(1+S_m)
\norm{U}_{C^{2m+4}(I_*)}^2$.
The actual positive scales satisfy $Y_{C,d}\gg_I D^{-C_0}$, so all output
frequencies and divisor factors are polynomially bounded even when $Y_{C,d}<1$.
Moreover $S_m$ is finite without induction: triangle and lattice counting give
$\E/\Sigma'\ll_{I'}\mathcal H'/F'\le\mathcal H$ for normalized tests.
An empty $S_m$ means there is no retained nonzero block.
Finally apply \eqref{eq:transfer-first-bound} with $j=2m+4$ and allocate
loss $\epsilon/4$ to each stage. Its order is
$2(2m+4)+4=4m+12$, and its total loss is at most $D^{\epsilon/2}$,
within the asserted $D^\epsilon$ budget.
\end{proof}

\subsection{Cube removal and the terminating induction}
For $\Sigma=XF$ and $\mathcal H\le\Sigma$, put
$H_c^3=\min(X,X^2/\mathcal H^2)$ and $L_b=X/\N(b)^3$.
Complete multiplicativity, including zeros, gives the exact cube inversion
\[
 X^{-1/2}\sum_n^*a_\xi(n)\chi_n(k)\chi_n(f)^4W(\N(n)/X)
 =\sum_{(h,S)=1}\frac{\mu(h)\overline{\alpha(h)}^{\,3}\Psi_k(h)^3}
 {\N(h)}T(X/\N(h)^3;k,f).
\]
Indeed the coefficient of the total cube index $b$ contains
$\sum_{h\mid b}\mu(h)=\mathbf1_{b=1}$. No coprimality with the
squarefree column is imposed. For $\N(h)\le H_c$, weighted Cauchy and the
completed interface with range $C_0+1$ give
$D^{\epsilon/2}\norm{W}_{C^{J_c}(I)}^2
(\mathcal H+\mathcal H^2FH_c^3/X)\ll
D^{\epsilon/2}\Sigma\norm{W}_{C^{J_c}(I)}^2$;
here $X/\N(h)^3\ge1$ and reciprocal-norm sums are logarithmic.
The larger range covers $\N(f)<2F$. If $H_c<1$, this short sum is empty.

Expanding the remaining $T$ and writing $b=hc$ yields the coefficient
\[
 \beta(b,f)=\overline{\alpha(b)}^{\,3}\xi(b)^3\mathbf1_{(b,f)=1}
       \sum_{\substack{h\mid b\\\N(h)>H_c}}\mu(h),
 \qquad |\beta(b,f)|\le\tau_{\rm div}(b).
\]
It multiplies the normalized original column at scale $L_b$ with weight
$\chi_b(k)^3/\N(b)$; no new column coefficient sequence is introduced.
Support requires $\N(b)^3\le vX$. Weighted Cauchy costs only the square
of $\sum_{\N(b)\le(vX)^{1/3}}\tau_{\rm div}(b)/\N(b)
\ll_I\log^2(2+X)$.
For $L_b\le1$ columns are empty unless $L_b\ge1/v$; otherwise direct
counting gives $\E(\mathcal H,L_b,F)\ll_I
(\mathcal H/L_b)\norm{W}_\infty^2\ll_I\Sigma\norm{W}_\infty^2$.
Thus, with a finite $J_c=J_c(\epsilon,C_0)$ independent of $I$,
\begin{equation}\label{eq:consolidated-cubes}
 \E(\mathcal H,X,F;\xi,W)\ll D^\epsilon
 \left(\Sigma\norm{W}_{C^{J_c}(I)}^2+
 \sup_{\N(b)>H_c,\ L_b>1}\E(\mathcal H,L_b,F;\xi,W)\right).
\end{equation}
If $\mathcal H^2\le X$, the supremum is empty. If $H_c<1$, the unit cube
belongs to the long range: its full divisor sum annihilates $b\ne1$,
but the larger displayed supremum is still a valid upper bound.

\begin{proposition}[Canonical estimate under a positive gap]
\label{prop:canonical-audit}
For $D\ge2$, $\kappa>0$, $C_0\ge1$, suppose $\mathcal H,X,F\ge1$,
$\Sigma=XF\le D^{C_0}$ and $\mathcal H\le\Sigma D^{-\kappa}$.
For every $\epsilon>0$ there is a finite integer $J=J(\kappa,C_0,\epsilon)$,
independent of $I$, such that
$\E(\mathcal H,X,F;\xi,W)\ll
D^\epsilon\Sigma\norm{W}_{C^J(I)}^2$.
\end{proposition}
\begin{proof}
Induct with the additional condition $\mathcal H\le D^{j\kappa}$.
For $j=0$, counting gives $\E\ll_I\mathcal HX\norm{W}_\infty^2
\le\Sigma\norm{W}_\infty^2$, so $J_0=1$ suffices.
For the step, take $m=J_j(\epsilon/3)$ and apply
\eqref{eq:consolidated-cubes} with loss $\epsilon/3$.
Every nonempty tail has $\mathcal H^2>X$ and $H_c^3=X^2/\mathcal H^2$.
Apply Proposition~\ref{prop:consolidated-transfer} at $L=L_b$, retaining
the independent target $\Sigma=XF$. Every returned child satisfies
\[
 \mathcal H'\le\frac{\mathcal HL_b}{\Sigma F}
 =\frac{\mathcal H}{\N(b)^3F^2}
 <\mathcal H\left(\frac{\mathcal H}{\Sigma}\right)^2
 \le D^{-2\kappa}\mathcal H,\qquad
 \frac{\mathcal H'}{\Sigma'}\le D^{-\kappa},\quad\Sigma'\le D^{C_0}.
\]
For $\mathcal H\le D^{(j+1)\kappa}$ it is covered by level $j$;
at the first step the sharper bound is $\mathcal H'<1$, so no child exists.
The strict contraction also holds for the unit cube when $H_c<1$.
Apply induction on $I'$, with loss $\epsilon/3$, and transfer with loss
$\epsilon/3$. Together with cube removal this spends exactly the allowed
$\epsilon$. It suffices to choose
$J_{j+1}(\epsilon)\ge\max\{J_c(\epsilon/3,C_0),4J_j(\epsilon/3)+12\}$.
Constants may depend on the finitely enlarged support interval; derivative
orders do not. Finally $j=\lceil C_0/\kappa\rceil$ covers
$\mathcal H\le\Sigma\le D^{C_0}$. All stages and derivative orders are finite.
\end{proof}

For later original-family applications, $\mathcal H<1$ gives an empty
nonzero row sum. If $X<1$, nonempty support forces $X\ge1/v$ and direct
counting gives $\E\ll_I\mathcal H\norm{W}_\infty^2$; a positive gap bounds
this by $\Sigma\norm{W}_\infty^2$. Exclusion shifts preserve that gap and
product, so the same short-scale argument applies after every shift.
These elementary endpoints and the finite descent add no new theta or
sieve cancellation premise to the declared interfaces.


\section{Returning to the original mean square}\label{sec:moment-final}

We now match the canonical estimate to
Source interface~\ref{in:initial}, including every boundary range.
The original row length is \(H=D^{1+\vartheta}\), whereas the
dual row length is \(\mathcal H\). From
\eqref{eq:initial-ranges}, with \(\Sigma=XF\), one has
\begin{equation}\label{eq:initial-gap-final}
 \mathcal H/\Sigma\le CD^{-\vartheta}/B\le CD^{-\vartheta}.
\end{equation}
For sufficiently large \(D\), this is at most
\(D^{-\kappa}\), with the fixed positive choice
\(\kappa=\vartheta/2\). The constant \(C\) changes the
starting threshold; it does not eliminate the exponent gap.

If the dual energy is nonempty, the compact support \(I=[a,b]\)
forces \(X\ge1/b\). Hence \(F=\Sigma/X\le bD/B\).
For \(X\ge1\), the canonical hypotheses hold with \(C_0=2\)
and the stated gap: \(\Sigma=D/B\le D\).
An exclusion shift \(d\mid\mathfrak b\), with \(\mathfrak b\) the
ideal exclusion index in the initial reduction, preserves
\(\Sigma,\mathcal H\) and the initial ranges. Its column length
can fall below one, so the next cases are needed before applying
the canonical proposition.

If \(\mathcal H<1\), the nonzero row set is empty. If \(X<1/b\),
the column set is empty. For \(1/b\le X<1\), the inner column sum
has bounded cardinality depending only on \(I\), and all coefficient
and character magnitudes are at most one. Since the \(f\)-shell has
\(O(F)\) ideals, direct counting gives
\[
 \E(\mathcal H,X,F;\xi,W)\ll_I
       \mathcal H\norm{W}_\infty^2
 \ll_I\Sigma\norm{W}_\infty^2
\]
for large \(D\) in the initial gap. The same reasoning applies
after every exclusion shift. The common-test interval of the
initial reduction is actually \(I_*=[a/8,2b]\);
there one uses the lower support bound \(X/\N(d)\ge1/(2b)\)
and constants for \(I_*\). Its enlargement is fixed.

For bounded \(D\ge2\), nonempty rows imply
\(B^2\le CD^{1-\vartheta}\), and nonempty columns imply
\(F\le bD/B\). Thus \(\Sigma\) is bounded away from zero and
all relevant ranges are bounded with permitted constants.
Counting covers them. These cases establish the auxiliary target
for every initial pair \(B,F\), not just the convenient long-column
subrange.

\begin{proposition}[Original moment under the audited interfaces]
\label{prop:original-moment}
The canonical estimate and Source interface~\ref{in:initial}
give \eqref{eq:main-moment}, with finite order
\(k=k(\vartheta,\epsilon)\) independent of \(I\).
\end{proposition}
\begin{proof}
For a common test \(U\) supported in \(I_*\), define
\[
 S_U(\mathfrak b,f,k)=
 \sfstar_{(n,\mathfrak b)=1}a_\xi(n)\chi_n(k)\chi_n(f)^4
                              U(\N(n)/X).
\]
In the following column moment, \(f\) ranges over squarefree primary
ideals prime to \(S\) with \(F\le\N(f)<2F\), and \(k\) over
nonzero elements with \(0<\N(k)\le\mathcal H\). Its initial
normalization is
\[
 \sum_{f,k}|S_U(\mathfrak b,f,k)|^2=XF\,\E_{\mathfrak b}.
\]
Apply the auxiliary target with loss \(\epsilon/4\).
Lemma~\ref{lem:exclusions} introduces the divisor count squared
of \(\mathfrak b\), where \(\N(\mathfrak b)\ll_I D\). Reserve a further
\(\epsilon/4\) for this loss, obtaining
\[
 \sum_{f,k}|S_U(\mathfrak b,f,k)|^2
 \ll D^{\epsilon/2}(XF)^2\norm{U}_{C^J(I_*)}^2.
\]
One factor \(XF\) is the auxiliary bound; the other undoes
the defining energy normalization. The coupled smooth-kernel
lemma in the initial reduction then uses \(q=2J+4\).
Its kernel norm is \(O_{I,q}(\norm{W}_{C^q(I)}^2)\),
so one dyadic \(B,F\) block has size at most
\[
 D^{\epsilon/2}\norm{W}_{C^q(I)}^2
  \frac{HB}{D^2}\cdot O(B)\cdot(XF)^2
 \ll D^{\epsilon/2}\norm{W}_{C^q(I)}^2H.
\]
Here \(XF=D/B\). The factors are the original Poisson
coefficient, the ideal count in the \(\mathfrak b\)-shell, and the
unnormalized column energy. The \(O_I((1+\log D)^2)\)
nonempty dyadic blocks use the remaining loss
\(\epsilon/2\). Sum the fixed finite ray-character expansion
and include \(Z\ll_I H\norm{W}_\infty^2\).
Consequently \(\mathcal M_D\ll
D^\epsilon H\norm{W}_{C^q(I)}^2\).

By \(\Phi\ge1\) on \([0,1]\) and \eqref{eq:MD},
\[
 \sum_{0<\N(u)\le H}|A_u(D)|^2
 \le D\mathcal M_D
 \ll D^\epsilon DH\norm{W}_{C^q(I)}^2
 =D^{2+\vartheta+\epsilon}\norm{W}_{C^q(I)}^2 .
\]
The lattice point \(u=0\) only has the unit column under
the source's character extension. That term vanishes for
large \(D\), and is counted for bounded \(D\). The order
\(J\) is chosen through finitely many canonical stages for
the fixed \(\vartheta,\epsilon\); its finite composition
with \(2J+4\) is independent of \(I\). This proves the claim.
\end{proof}



% Prepared for Edward Baker, 12 October 2026, with substantial LLM assistance.
% Model: GPT-6 (Codex), inherited configuration. Exact serving variant and
% configured reasoning effort are unavailable and are not inferred.
% Fragment only: no preamble, author field, bibliography, or compilation.
% Source: openai/math fd4aeeb2ee4fc729c18d98444fed42fd0529eeeb, paper2.tex.
% SHA-256: d9a8f15aa770cf883d0eabd2b775fad694ce20b44cba7928f5c0c9a6d8750d4d.

\section{Fixed-character extraction and the zero-free implication}
\label{sec:final-extraction}

This section records the consequences of the original moment assembled
in the preceding sections. Its conclusions are conditional on that
moment estimate and on the classical continuation and prime-counting
inputs stated below. The argument reorganizes the concluding deductions
of \cite{Source}; it does not independently establish the new arithmetic
mean square or furnish a separate proof of the quasi-Riemann hypothesis.

Let \(K=\mathbb Q(\sqrt{-3})\), and retain its ring of integers
\(\OO\), the ideal norm \(\N\), and the zero-extended sextic characters
\(\chi_n\). Ideals prime to the primes above 3 are represented by their
unique primary generators. Fix a finite-order Hecke character \(\nu\)
and a finite set \(S\) containing its conductor primes and the primes
above 2 and 3. For a smooth complex-valued test \(W\) of compact support
in \((0,\infty)\), put
\begin{equation}\label{eq:final-A}
 A_u(D)=\sum_{(n,S)=1}^{*}\mu(n)\nu(n)\chi_n(u)
                  W\!\left(\frac{\N(n)}D\right),\qquad D>0,
\end{equation}
where the star means squarefree integral ideals, with their primary
representatives. The fixed-character moment input is the following:
for every \(0<\vartheta\le1/10\) and \(\eta>0\), there is an integer
\(k=k(\vartheta,\eta)\ge1\), independent of the character, its fixed
excluded set and the support interval, such that for every compact
interval \(I\subset(0,\infty)\), every smooth \(W\) supported in
\(I\), and \(D\ge2\),
\begin{equation}\label{eq:final-moment-input}
 \sum_{0<\N(u)\le D^{1+\vartheta}}|A_u(D)|^2
 \ll_{\nu,S,I,\vartheta,\eta}
       \norm{W}_{C^k(I)}^2 D^{2+\vartheta+\eta},
 \qquad
 \norm{W}_{C^k(I)}=\max_{0\le j\le k}\sup_{y\in I}|W^{(j)}(y)|.
\end{equation}
The sum over \(u\) is over nonzero elements of \(\OO\). In particular,
the input includes the principal character \(\nu=1\).

\subsection{Coherent sixth-power rows}

\begin{proposition}[Extraction from the original moment]
\label{prop:final-extraction}
Assume \eqref{eq:final-moment-input}. For every \(\delta>0\) there is
an integer \(k_\delta\), depending only on \(\delta\), such that,
for every fixed \(\nu,S,I\) as above and every smooth \(W\) supported
in \(I\),
\begin{equation}\label{eq:final-saving}
 |A_1(D)|\ll_{\nu,S,I,\delta}
          \norm{W}_{C^{k_\delta}(I)}D^{11/12+\delta}
          \qquad(D\ge2).
\end{equation}
\end{proposition}

\begin{proof}
First fix \(0<\vartheta\le1/10\), and set
\[
 H=D^{1+\vartheta},\qquad Y=H^{1/6}.
\]
For each prime ideal \(\mathfrak p\notin S\) with
\(Y/2<\N(\mathfrak p)\le Y\), let \(p\) denote its primary
generator. The zero-extension convention gives the exact identity
\begin{equation}\label{eq:final-coherent}
 \chi_n(p^6)=\mathbf 1_{\mathfrak p\nmid n}.
\end{equation}
Indeed, a nonzero sextic character value has sixth power one; if
\(\mathfrak p\mid n\), its value is zero. Moreover
\(\N(p^6)\le H\). These are distinct element rows: an equality
\(p_1^6=p_2^6\) implies equality of the sixth powers of their prime
ideals, hence equality of the ideals and of their primary generators.

Write \(I=[a,b]\). The Eisenstein lattice gives
\(\#\{\mathfrak a:\N(\mathfrak a)\le R\}\ll1+R\).
Consequently, uniformly over the selected primes,
\[
 \#\{n:\mathfrak p\mid n,\ aD\le\N(n)\le bD\}
 \le\#\{\mathfrak a:\N(\mathfrak a)\le bD/\N(\mathfrak p)\}
 \ll_I D/Y.
\]
Here \(D/Y=D^{(5-\vartheta)/6}\ge1\), so the constant term in the
lattice estimate is absorbed. Since \(|\mu(n)\nu(n)|\le1\),
\begin{equation}\label{eq:final-prime-error}
 |A_1(D)-A_{p^6}(D)|\ll_I\norm{W}_\infty D/Y.
\end{equation}
No cancellation or varying-ray-class prime count is used in this error.

We import the classical fixed-field prime ideal theorem, as in the
source argument \cite{Source}; see also the primary presentation
\cite{Burgin}. With the character and \(S\) fixed
before \(D\) grows, it gives
\[
 J=\#\{\mathfrak p\notin S:Y/2<\N(\mathfrak p)\le Y\}
       \sim\frac{Y}{2\log Y}.
\]
Removing the finite set \(S\) changes the starting threshold, not this
asymptotic. Positivity, \eqref{eq:final-prime-error}, and the moment
with loss \(\eta\) now imply
\begin{align*}
 |A_1(D)|^2
 &\le\frac2J\sum_p|A_{p^6}(D)|^2
       +O_I\!\left(\norm{W}_\infty^2D^2/Y^2\right)\\
 &\ll_{\nu,S,I,\vartheta,\eta,\lambda}
       \norm{W}_{C^{k(\vartheta,\eta)}(I)}^2
       D^{11/6+5\vartheta/6+\eta+\lambda}
       +\norm{W}_\infty^2D^{5/3-\vartheta/3},
\end{align*}
where \(\lambda>0\) absorbs \(\log Y\ll_\lambda D^\lambda\).
Thus the two exponents after taking square roots are
\[
 \frac{11}{12}+\frac{5\vartheta}{12}
                  +\frac{\eta+\lambda}{2},
 \qquad \frac56-\frac{\vartheta}{6}.
\]
For the requested \(\delta>0\), choose in order
\[
 \vartheta=\min\{1/10,3\delta/5\}>0,
 \qquad \eta=\lambda=\delta/2,
 \qquad k_\delta=k(\vartheta,\eta).
\]
The loss in the first exponent is at most \(3\delta/4<\delta\),
and the second exponent is below \(11/12\). The permitted constant
also covers the bounded initial range before the prime asymptotic
applies. This proves \eqref{eq:final-saving}. The proof uses only the
unit-modulus bound on \(\nu\), so it applies unchanged to \(\nu=1\).
\end{proof}

\subsection{Mellin holomorphy and a test at a fixed alleged zero}

Let \(L_K(s,\nu)\) denote the Hecke \(L\)-function with its actual
finite conductor, and let \(L_K^S\) be its Euler product with primes
in \(S\) removed. At a conductor prime use \(\nu(\mathfrak p)=0\).
Then
\begin{equation}\label{eq:final-missing-factors}
 L_K^S(s,\nu)=L_K(s,\nu)
       \prod_{\mathfrak p\in S}
       \bigl(1-\nu(\mathfrak p)\N(\mathfrak p)^{-s}\bigr).
\end{equation}
Every factor in this finite product, and its inverse, is holomorphic and
nonvanishing on \(\Re s>0\), because
\(|\nu(\mathfrak p)\N(\mathfrak p)^{-s}|<1\) there whenever
the character value is nonzero.

We import classical Hecke continuation in the form supplied, for
example, by \cite[Theorem 5.22 and Section 5.11]{Poonen}. A finite-order
character at the complex place is trivial, since a continuous map from
the connected group \(\mathbb C^\times\) into a finite group is
constant. Thus the finite Hecke \(L\)-function is holomorphic in
\(\Re s>0\), except for the possible simple pole at \(s=1\) of the
principal character. This continuation theorem is an external input.

\begin{proposition}[Conditional Hecke zero-free conclusion]
\label{prop:final-hecke}
Assume \eqref{eq:final-moment-input} and the classical inputs just
stated. For every fixed finite-order Hecke character \(\nu\) of
\(K\), \(L_K(s,\nu)\) has no zeros in
\(\Re s>11/12\). The principal pole at \(s=1\) is permitted.
\end{proposition}

\begin{proof}
For a fixed test \(W\), define
\[
 \widehat W(s)=\int_0^\infty W(y)y^s\,\frac{dy}{y},
 \qquad
 M_W(s)=\int_0^\infty A_1(D)D^{-s}\,\frac{dD}{D}.
\]
If \(\operatorname{supp}W\subset[a,b]\), then \(A_1(D)=0\)
for \(D<1/b\), because every nonzero integral ideal has norm at
least one. On any bounded positive interval in \(D\), only finitely
many ideals occur. Its contribution to \(M_W\) is therefore entire.

Put \(\alpha=11/12\), and let \(T\) be a compact subset of
\(\{\Re s>\alpha\}\). Choose one positive loss
\(\delta<\min_{s\in T}\Re s-\alpha\) in
\eqref{eq:final-saving}. The resulting tail majorant is integrable
with respect to \(dD/D\), uniformly for \(s\in T\); every derivative
in \(s\) adds an integrable power of \(\log D\). Consequently
\(M_W\) is holomorphic on all of \(\Re s>\alpha\). A separate loss
on each compact subset suffices, with \(W\) unchanged.

For \(\sigma=\Re s>1\), absolute interchange is justified by
\[
 \sum_n\int_0^\infty
   |\mu(n)\nu(n)W(\N(n)/D)|D^{-\sigma}\,\frac{dD}{D}
 \le
 \left(\int_0^\infty|W(y)|y^\sigma\,\frac{dy}{y}\right)
 \sum_n\N(n)^{-\sigma}<\infty.
\]
Changing variables \(y=\N(n)/D\) gives
\begin{equation}\label{eq:final-mellin-identity}
 M_W(s)=\widehat W(s)
      \sum_{(n,S)=1}\mu(n)\nu(n)\N(n)^{-s}
      =\frac{\widehat W(s)}{L_K^S(s,\nu)}
      \qquad(\Re s>1).
\end{equation}
The meromorphic identity theorem therefore extends
\begin{equation}\label{eq:final-product-identity}
 L_K^S(s,\nu)M_W(s)=\widehat W(s)
       \qquad(\Re s>\alpha).
\end{equation}
This formulation does not assume a holomorphic reciprocal at an
unknown zero. In the principal case it makes the product's possible
singularity at 1 removable; its nonzero residue also forces
\(M_W(1)=0\). That pole is not a zero.

Suppose that \(L_K(s,\nu)\) has a zero \(\rho\) with
\(\Re\rho>\alpha\). Choose a nonzero nonnegative
\(\phi\in C_c^\infty((1,2))\), and fix
\[
 W_\rho(y)=\exp(-\rho\log y)\phi(y),
\]
using the real logarithm on positive \(y\). This is an admissible
complex smooth test, fixed after \(\rho\) has been specified, and
\[
 \widehat W_\rho(\rho)=\int_0^\infty\phi(y)\,\frac{dy}{y}>0.
\]
Its derivative norms may depend on \(\rho\); the fixed-test estimate
allows that dependence. The finite factors in
\eqref{eq:final-missing-factors} preserve the zero at \(\rho\).
At that point \(L_K^S\) is regular, while \(M_{W_\rho}\) is
holomorphic. Evaluating \eqref{eq:final-product-identity} gives
\(0=\widehat W_\rho(\rho)>0\), a contradiction. This also handles
a zero at 1 for a nonprincipal character.
\end{proof}

\subsection{Dirichlet characters and the norm lift}

Let \(\chi\) be a Dirichlet character modulo \(q\), and put
\(\Delta=\chi_{-3}\). Define, initially on integral ideals prime
to \(q\OO\),
\[
 \nu_\chi(\mathfrak a)=\chi\bigl(\N(\mathfrak a)\bigr).
\]
Extend it multiplicatively to fractional ideals prime to that modulus,
using \(\chi\) on positive rational \(q\)-units. For a ray principal
ideal generated by \(\beta\equiv1\pmod{q\OO}\),
\(\N(\beta)\equiv1\pmod q\), so \(\nu_\chi((\beta))=1\).
The norm lift is therefore a finite-order ray character with actual
finite conductor \(\mathfrak f_\chi\mid q\OO\) and trivial infinity
type. Oddness of \(\chi\) introduces no extra infinite type: the
complex norm is positive, and the real sign character is trivial on
its image.

Write \(L_K^{[q]}(s,\nu_\chi)\) for the Euler product omitting all
prime ideals above primes dividing \(q\). The local comparison can
be made directly. For \(p\nmid q\), put \(z=\chi(p)p^{-s}\).
The factors are
\[
\begin{array}{c|c|c}
 \text{rational prime} & L_K^{[q]}(s,\nu_\chi)
                        & L(s,\chi)L(s,\chi\Delta)\\ \hline
 p\equiv1\pmod3 &(1-z)^{-2}&(1-z)^{-1}(1-z)^{-1}\\
 p\equiv2\pmod3 &(1-z^2)^{-1}&(1-z)^{-1}(1+z)^{-1}\\
 p=3, 3\nmid q &(1-\chi(3)3^{-s})^{-1}
                    &(1-\chi(3)3^{-s})^{-1}.
\end{array}
\]
The first row comes from two prime ideals of norm \(p\), the second
from one prime ideal of norm \(p^2\), and the last from the ramified
prime of norm 3, with \(\Delta(3)=0\). The inert row includes
\(p=2\); its exclusion during the arithmetic argument does not omit
it from this full Euler comparison. At \(p\mid q\), both sides of
the modulus-product comparison have local factor one. Here
\(\chi\Delta\) is regarded as the pointwise product character with
natural modulus \(\operatorname{lcm}(q,3)\). Hence, first for
\(\Re s>1\) and then by meromorphic continuation,
\begin{equation}\label{eq:final-normlift}
 L_K(s,\nu_\chi)=E_q(s)L(s,\chi)L(s,\chi\Delta),
\end{equation}
where the actual-conductor correction is
\[
 E_q(s)=\prod_{\substack{\mathfrak p\mid q\OO\\
                         \mathfrak p\nmid\mathfrak f_\chi}}
       \bigl(1-\nu_\chi(\mathfrak p)\N(\mathfrak p)^{-s}\bigr)^{-1}.
\]
This finite product and its inverse are holomorphic and nonvanishing
on \(\Re s>0\).

The conductor distinction is necessary. The norm lift of \(\Delta\)
is trivial away from 3, since split prime norms are 1 modulo 3 and inert
prime norms are squares. Its actual conductor is one, and
\[
 L_K(s,\nu_\Delta)=\zeta_K(s),\qquad
 E_3(s)=(1-3^{-s})^{-1}.
\]
Thus removing the finite correction in
\eqref{eq:final-normlift} would give the wrong exact identity, even
though the positive-half-plane zero comparison is unaffected.

\begin{proposition}[Conditional Dirichlet conclusion]
\label{prop:final-dirichlet}
Under the assumptions of Proposition~\ref{prop:final-hecke}, every
Dirichlet \(L\)-function, including those of imprimitive and principal
characters, has no zeros in \(\Re s>11/12\). Principal poles at 1
are permitted.
\end{proposition}

\begin{proof}
If \(\chi\) is induced by a primitive character \(\chi^*\) of
conductor \(q_0\mid q\), then comparison of Euler products gives
\[
 L(s,\chi)=L(s,\chi^*)
       \prod_{\substack{p\mid q\\p\nmid q_0}}
       (1-\chi^*(p)p^{-s}).
\]
All finite factors are nonzero on \(\Re s>0\). The same statement
applies to the possibly imprimitive character \(\chi\Delta\).
Classical Dirichlet continuation \cite[Section 25.15]{DLMF} says
that nonprincipal primitive factors are entire and principal factors
have only a simple pole at 1. Thus, away from 1, a zero of either factor
in \eqref{eq:final-normlift} forces a zero of the Hecke function,
contrary to Proposition~\ref{prop:final-hecke}.

At 1, if one primitive factor is principal, the other primitive factor
is \(\Delta\): either \(\chi^*=1\), or the primitive character
inducing \(\chi\Delta\) is principal and \(\chi^*=\Delta\).
They cannot both be principal. The quadratic value is nonzero, as is
also seen directly from the convergent positive paired series
\[
 L(1,\Delta)=\sum_{j\ge0}
       \left(\frac1{3j+1}-\frac1{3j+2}\right)>0.
\]
Nonzero imprimitive Euler corrections cannot change this fact, so the
principal pole cannot cancel a zero of the other factor. If neither
primitive factor is principal, both are holomorphic at 1, and their
product cannot vanish there by the Hecke conclusion. These cases cover
all characters.
\end{proof}

\begin{remark}[Scope of the conclusion]
The boundary \(11/12\) is common to all fixed characters. The moment
and extraction derivative orders depend on the chosen exponent losses,
but not on the fixed character or support interval. Constants,
prime-counting thresholds, and the test norms used for a fixed alleged
zero may depend on that character, its fixed conductor and excluded
set, the support, and the alleged zero. The argument establishes a
qualitative assertion for every fixed character and every fixed zero;
it does not supply quantitative uniformity over conductor or height,
nor the source's further quantitative corollaries. All conclusions
here retain their dependence on the preceding original moment and the
stated classical imports.
\end{remark}


\section{Finite diagnostics and reproducibility}\label{sec:verification}

% Prepared for Edward Baker, 12 October 2026, with substantial LLM assistance.
% Model: GPT-6 (Codex), inherited configuration. Exact serving variant and
% configured reasoning effort are not exposed and are not inferred.
% Internal record inventory; no new tests were run for this writing task.
% Requires booktabs and tabularx. No additional manuscript macros are needed.
% The table is an inventory of saved records, not a combined validation count.

The saved diagnostics provide finite checks of the arithmetic identities
and parameter bookkeeping used in the audit. Table~\ref{tab:finite-audit}
records the counters in the verified run records. These counters have
different meanings, include repeated checks within enumerations, and
are not added together or interpreted as numbers of independent theorem
obligations. The first branch reports cases rather than a total assertion
count; the theta row is a replay of an existing suite, which also contains
primitive-Poisson tests.

\begin{table}[htbp]
\centering
\small
\setlength{\tabcolsep}{4pt}
\renewcommand{\arraystretch}{1.12}
\begin{tabularx}{\linewidth}{@{}p{0.16\linewidth}p{0.29\linewidth}X@{}}
\toprule
Audit stage & Recorded diagnostics & Principal finite coverage \\
\midrule
Selected branch & 729 regrouping cases; 4,913 preimage normalizations;
one unsigned control
& 55 retained scale cases and 65 exclusion shifts; rational scale and
necessary-support inequalities. \\
Theta replay & 69,881 exact assertions, including eight counterchecks
& 4,860 cusp constructions, all three cusp cases, local finite
Fourier transforms and additive CRT. \\
Initial reduction & 113,748 exact assertions; 19,868 floating comparisons;
12 negative-control families
& 4,374 gcd-bijection inputs; finite ray phases and zero masks; direct
Gauss sums at six split and two inert primes, and 28 prime pairs. \\
Completed mean square & 198,323 exact assertions; 11 negative-control families
& 40 local reindexing sums, 7,776 Eisenstein rows, valuation budgets,
general-row zero extensions and inverse-square normalization. \\
Full transfer & 4,186,781 exact assertions; 11 negative-control families
& 62,208 signed regrouping cases and 438,012 Eisenstein preimages;
mask identities, scale changes and finite descent bookkeeping. \\
Final assembly & 622 exact assertions; 13 negative-control families
& Moment normalization, sixth-power extraction, positive loss budgets,
short-scale boundaries and local Euler polynomials. \\
\bottomrule
\end{tabularx}
\caption{Saved finite diagnostic records of 12 October 2026. The theta
suite was originally recorded on 9 October and replayed unchanged on
12 October. Negative controls alter a sign, mask, normalization or
admissibility rule and detect the resulting failure.}
\label{tab:finite-audit}
\end{table}

All suites except the initial reduction use exact arithmetic without
floating comparisons. Its ordinary floating diagnostics have tolerance
$10^{-10}$ and maximum recorded absolute discrepancy
$4.091524622677201\times10^{-15}$; no interval certification is supplied.
The prime-pair test sets the original ray twist to one, so its extension
to a general fixed twist rests on the algebraic ray-character argument.
The full-transfer record additionally lists 271 retained scale cases,
309 exclusion-shifted children and 1,311 eligible cube-tail parameter
cases. Eligibility is a statement about parameter inequalities, not
about nonzero M\"obius coefficients, energies or smooth kernels. No
actual smooth kernel or theta function is numerically evaluated by these
suites. The recorded hashes and Python versions fix the tested artifacts;
the legacy theta checker embeds a historical PDF hash, which the later
version comparison verified separately, rather than reading that PDF
during the replay.

These checks do not establish analytic Poisson summation, global cubic
reciprocity, theta automorphy and cusp-coefficient theorems, the quadratic
large sieve, infinite estimates, Hecke continuation or Dirichlet
nonvanishing. The audit compares the cusp formulas directly with the
Dunn--Radziwi\l\l\ v3 presentation. Patterson's original proof was
not independently read; the underlying cubic theta theorem remains an
import. The analytic conclusions in the manuscript are therefore stated
under the recorded interfaces and imports. They concern fixed characters
and a fixed excluded set; conductor or height uniformity is not
provided by these records. Internal parallel LLM read-throughs are not
independent mathematical review or formal verification.

% Record basenames, all under arithmetic_source_audit/numerics/:
% transfer_regrouping_record_20261012.json
% theta_version_crosswalk_record_20261012.json
% initial_arithmetic_record_20261012.json
% completed_bookkeeping_record_20261012.json
% full_transfer_bookkeeping_record_20261012.json
% final_assembly_record_20261012.json
% Theta legacy record, replayed verbatim:
% quasi-rh-character-amplification/numerics/
%   upstream_reflection_interfaces_record_20261009.json


\section{Conclusion of the audit}

The conditional assembly in Theorem~\ref{thm:assembly} now follows
from Propositions~\ref{prop:completed}, \ref{prop:consolidated-transfer}
and \ref{prop:canonical-audit}, Proposition~\ref{prop:original-moment},
and Propositions~\ref{prop:final-extraction},
\ref{prop:final-hecke}, and \ref{prop:final-dirichlet}.
The principal M\"obius estimate is extracted from the positive
original moment; it was not assumed as a separate cancellation
theorem. The special coefficient class and its zero extensions
survive the signed return, and the positive gap supplies
a strictly decreasing recursion parameter.

No further unassembled source link was identified in this
bounded main-theorem audit. The next level of scrutiny is
independent mathematical review of the full argument and its
deep imported interfaces, particularly the original cubic theta
theorem and the cusp coefficients used here. The quantitative
introduction corollaries, effective conductor or height estimates,
and later stronger strips were outside this assembly.

\section*{Acknowledgement and research provenance}
\addcontentsline{toc}{section}{Acknowledgement and research provenance}

This manuscript was drafted for Edward Baker with substantial
assistance from GPT-6 (Codex), including parallel internal
read-throughs, mathematical organization, and finite diagnostic
bookkeeping. The session used its inherited configuration; the
exact serving variant and configured reasoning effort were not
exposed and are not inferred. The model-assisted read-throughs
are internal checks, not independent specialist validation.
The source argument and its claimed theorem belong to the
authors of \cite{Source}; the audit's derivations and their
conditional scope are identified in the text.

\appendix
\section{Repository evidence and replay guide}\label{app:evidence}

The stable manuscript is
\path{arithmetic_source_audit.tex} in the investigation root.
Incremental research remains in \path{notes/}, assessments in
\path{reviews/}, and standard-library diagnostics and small
records in \path{numerics/}. This manuscript consolidates those
records rather than replacing their detailed source locators.

\begin{center}\small
\begin{tabularx}{\textwidth}{@{}lX@{}}
\toprule
Stage & Detailed local records\\
\midrule
Initial arithmetic and Poisson &
\path{notes/INITIAL_ARITHMETIC_CONVERSION_AUDIT_20261012.md};
\path{notes/INITIAL_POISSON_REDUCTION_AUDIT_20261012.md}\\
Theta inputs and normalization &
\path{notes/THETA_CUSP_AND_UNIFORMITY_AUDIT_20261012.md};
\path{notes/THETA_MELLIN_AND_KERNEL_AUDIT_20261012.md};
\path{notes/THETA_PRIMARY_SOURCE_COMPARISON_20261012.md}\\
Completed mean square &
\path{notes/COMPLETED_MEAN_SQUARE_LOCAL_LEDGER_20261012.md};
\path{notes/COMPLETED_MEAN_SQUARE_SIEVE_AND_SUMMATION_20261012.md};
\path{notes/COMPLETED_MEAN_SQUARE_GENERAL_ROWS_20261012.md}\\
Full transfer and termination &
\path{notes/TWO_POISSON_ARITHMETIC_LEDGER_20261012.md};
\path{notes/TWO_POISSON_SMOOTH_AND_SHORT_RANGES_20261012.md};
\path{notes/TWO_POISSON_RETURN_AND_TERMINATION_20261012.md}\\
Final deductions &
\path{notes/FINAL_ORIGINAL_MOMENT_ASSEMBLY_20261012.md};
\path{notes/FINAL_EXTRACTION_AND_BASE_CHANGE_20261012.md};
\path{notes/FINAL_MELLIN_AND_ZERO_FREE_ASSEMBLY_20261012.md}\\
Combined assessment &
\path{reviews/FINAL_ASSEMBLY_REVIEW_20261012.md}\\
\bottomrule
\end{tabularx}
\end{center}

From this folder, the final finite diagnostic can be replayed with
\begin{quote}\small\ttfamily
python3 numerics/check\_final\_assembly.py\\
\hspace*{1em}--output /tmp/final-assembly-replay.json
\end{quote}
To verify the source version as well, add
\path{--source /path/to/paper2.tex}, using the immutable URL and hash
in \path{notes/source_provenance_20261012.json}. A supplied
mismatched source is rejected before diagnostic arithmetic or
output; a run without a source explicitly records that its hash
was not verified. Full commands and limitations for every suite
are in \path{numerics/README.md}. The third-party source and PDFs
remain outside this repository, in accordance with its large-file
policy.

\begin{thebibliography}{9}
\bibitem{Source}
OpenAI, \emph{The Quasi-Riemann Hypothesis}, manuscript dated
October 5, 2026, openai/math revision
\texttt{fd4aeeb2ee4fc729c18d98444fed42fd0529eeeb},
\path{preprints/The-Quasi-Riemann-Hypothesis-October-5-2026/build/paper2.tex}.
\href{https://github.com/openai/math/blob/fd4aeeb2ee4fc729c18d98444fed42fd0529eeeb/preprints/The-Quasi-Riemann-Hypothesis-October-5-2026/build/paper2.tex}{Immutable source}.

\bibitem{GL}
L. Goldmakher and B. Louvel,
\emph{A quadratic large sieve inequality over number fields},
arXiv:1112.1642v2, Definition 1, Theorem 1.1 and Section 2.
\url{https://arxiv.org/html/1112.1642v2}.

\bibitem{DR}
A. Dunn and M. Radziwi{\l}{\l},
\emph{Bias in cubic Gauss sums: Patterson's conjecture},
arXiv:2109.07463v3, May 14, 2024, Section 5 and Appendix A.
\url{https://arxiv.org/html/2109.07463v3}.

\bibitem{Pat}
S. J. Patterson, \emph{A cubic analogue of the theta series},
Journal f\"ur die reine und angewandte Mathematik 296 (1977),
125--161. \url{https://doi.org/10.1515/crll.1977.296.125}.
The original proof was not independently accessed in this audit.

\bibitem{Poonen}
B. Poonen, \emph{Tate's thesis}, MIT 18.786 lecture notes,
Spring 2015, Theorem 5.22 and Section 5.11.
\url{https://math.mit.edu/~poonen/786/notes.pdf}.

\bibitem{DLMF}
NIST Digital Library of Mathematical Functions,
Section 25.15, \emph{Dirichlet \(L\)-functions}.
\url{https://dlmf.nist.gov/25.15}.

\bibitem{Burgin}
A. Burgin, \emph{A New Elementary Proof of Landau's Prime Ideal
Theorem, and Associated Results}, arXiv:2406.08565v1.
\url{https://arxiv.org/html/2406.08565v1}.
\end{thebibliography}
\end{document}
